\documentclass[12pt,a4paper]{article}
\usepackage[T1]{fontenc}
\usepackage[utf8]{inputenc}
\usepackage[polish]{babel}
\usepackage{mathptmx}
\usepackage{amsmath}
\usepackage[margin=25mm]{geometry}
\usepackage{microtype}
\usepackage{booktabs,tabularx}
\usepackage{graphicx}
\usepackage{listings}
\usepackage{caption}
\usepackage{float}
\usepackage[hidelinks]{hyperref}

\linespread{1.15}
\setlength{\parindent}{1.25cm}
\setlength{\parskip}{0pt}
\setlength{\emergencystretch}{2em}
\setlength{\tabcolsep}{9pt}
\renewcommand{\arraystretch}{1.1}
\captionsetup{font=small,labelfont=bf,labelsep=colon}
\captionsetup[table]{position=top,name=Tabela}
\captionsetup[figure]{name=Rysunek}
\lstset{language=C++,basicstyle=\ttfamily\small,columns=fullflexible,
  keepspaces=true,breaklines=true,showstringspaces=false,
  aboveskip=5pt,belowskip=5pt}

\newcommand{\wykres}[3]{%
  \begin{figure}[H]
  \centering
  \includegraphics[width=0.9\textwidth]{#1}
  \caption{#2}\label{#3}
  \end{figure}}

\hypersetup{pdftitle={Porównanie wydajności mnożenia macierzy na CPU i GPU},
  pdfsubject={Raport}}

\begin{document}
\hypersetup{pageanchor=false}
\begin{titlepage}
\centering
{\Large Politechnika Wrocławska\par}
\vspace{12mm}
{\large Akceleracja obliczeń w przetwarzaniu danych\par}
\vspace{20mm}
{\LARGE Porównanie wydajności mnożenia macierzy\\[3mm]
na CPU i GPU\par}
\vspace{8mm}
{\large Raport\par}
\vspace{18mm}
\begin{tabular}{c}
Michał Musielak 281092\\[3mm]
Klaudyna Kijek 280891\\[3mm]
Mikołaj Brzuśnian 280895\\[3mm]
Yuliia Nudha 282876
\end{tabular}
\vfill
\begin{tabular}{rl}
Prowadzący: & Michał Zając
\end{tabular}
\vspace{12mm}
\end{titlepage}
\hypersetup{pageanchor=true}

\tableofcontents
\clearpage
\raggedbottom

\section{Cel projektu}

Celem projektu było porównanie czasu klasycznego mnożenia macierzy
w trzech wersjach: sekwencyjnej na CPU, równoległej na CPU i równoległej
na GPU. Zbadano wpływ rozmiaru macierzy i liczby wątków na wydajność
oraz wykorzystano profiler do analizy pracy GPU.

\section{Przyjęte założenia i środowisko testowe}

Program napisano w C++17 z użyciem \texttt{std::thread} i CUDA 13.4.
Kompilację wykonano z optymalizacją \texttt{-O3}.
Macierze kwadratowe zapisano wierszami, używając typu \texttt{float}.
Wszystkie wersje realizują ten sam algorytm i otrzymują te same dane.
Nie korzystano z gotowych bibliotek do mnożenia macierzy.

\begin{table}[H]
\caption{Środowisko testowe}
\label{tab:sprzet}
\centering
\begin{tabularx}{\textwidth}{@{}lX@{}}
\toprule
Element & Specyfikacja\\
\midrule
CPU & AMD Ryzen 7 3700X, 8 rdzeni, 16 wątków\\
GPU & NVIDIA GeForce RTX 3060 Ti, 8 GiB\\
RAM & 31,26 GiB\\
System & CachyOS Linux\\
\bottomrule
\end{tabularx}
\end{table}

Na czas badań ograniczono liczbę działających aplikacji.
Zegar GPU ustawiono na 1500\,MHz, a pamięci na 6801\,MHz
opcjami \texttt{-lgc} i \texttt{-lmc} narzędzia
\texttt{nvidia-smi} \cite{nvsmi}. Monitoring potwierdził te wartości.
Dla CPU wybrano tryb \texttt{performance}, zadano 3,60\,GHz
i wyłączono boost. Odczyty CPU wynosiły 3,36--3,61\,GHz;
żądana częstotliwość nie gwarantuje idealnie stałego taktowania
\cite{cpufreq}. Program mógł korzystać z ośmiu rdzeni fizycznych.
Po badaniach przywrócono poprzednie ustawienia komputera.

\section{Wstęp teoretyczny}

Iloczyn macierzy $A$ i $B$ o rozmiarze $N\times N$ wyznacza wzór
\begin{equation}
c_{ij}=\sum_{k=0}^{N-1} a_{ik}b_{kj}.
\end{equation}
Każdy element wyniku jest sumą iloczynów elementów wiersza $A$
i kolumny $B$. Przykładowo:
\[
\begin{bmatrix}1&2\\3&4\end{bmatrix}
\begin{bmatrix}5&6\\7&8\end{bmatrix}
=\begin{bmatrix}19&22\\43&50\end{bmatrix}.
\]
Mnożenie macierzy wykorzystuje się m.in. w grafice i obliczeniach
naukowych. Klasyczny algorytm ma złożoność $O(N^3)$.

Elementy wyniku są niezależne, dlatego można obliczać je równolegle.
CPU dzieli pracę między kilka rdzeni, a GPU wykonuje wiele podobnych
zadań jednocześnie. Korzyść z równoległości ograniczają koszty
obsługi wątków, dostępu do pamięci oraz przesyłania danych do GPU.
Przy małych macierzach koszty te mogą przewyższać oszczędność czasu
samych obliczeń.

\section{Implementacja}

\subsection{Trzy wersje algorytmu}

Wersja sekwencyjna wykonuje trzy pętle: po wierszach, kolumnach
i składnikach sumy. Wersja równoległa CPU dzieli wiersze wyniku
między wątki \texttt{std::thread}. Wątki zapisują rozłączne części
macierzy, więc nie wymagają blokad. Przed odebraniem wyniku program
czeka na zakończenie wszystkich wątków.

W wersji CUDA jeden wątek oblicza jeden element wyniku.
Wątki uruchamiane są w blokach po 256. Poniżej przedstawiono
kod jądra GPU:
\begin{lstlisting}
__global__ void gpu_kernel(const float* a,
    const float* b, float* c, int n) {
    int index = blockIdx.x * blockDim.x
              + threadIdx.x;
    if (index >= n*n) return;
    int i = index / n;
    int j = index % n;
    float sum = 0;
    for (int k = 0; k < n; ++k)
        sum += a[i*n+k] * b[k*n+j];
    c[index] = sum;
}
\end{lstlisting}
Przed obliczeniami dane są kopiowane do GPU, a po nich wynik wraca
do pamięci CPU. Ważne było poprawne obsłużenie niepełnego bloku GPU,
nierównego podziału wierszy CPU oraz zakończenia pracy GPU przed
zatrzymaniem pomiaru czasu.

\subsection{Sprawdzanie poprawności}

Sprawdzono znany iloczyn $2\times2$, macierze zerowe, jednostkowe,
wypełnione jedynkami oraz dane losowe, ułamkowe i ujemne.
Testy obejmowały również rozmiary 17 i 33, które nie dzielą się
równo między wszystkie badane wątki. Wszystkie dziewięć przypadków
przeszło testy w każdej wersji.

Dopuszczalną różnicę zmiennoprzecinkową określono warunkiem
\begin{equation}
|c-c_{\mathrm{ref}}|\leq10^{-3}+10^{-4}|c_{\mathrm{ref}}|.
\end{equation}
Dodatkowe porównanie z niezależnym wynikiem typu \texttt{double}
objęło 108 sprawdzeń dla rozmiarów od 1 do 513. Wszystkie były
poprawne; największy błąd bezwzględny wyniósł
$3{,}52\cdot10^{-5}$.

\section{Metodologia przeprowadzania badań}

Macierze wypełniano liczbami pseudolosowymi z przedziału $[-1,1]$,
ze stałym ziarnem 42. Zbadano rozmiary $N=64$, 128, 256 i 512
oraz CPU równoległy z 1, 2, 4 i 8 wątkami.
Wykonano trzy pełne serie. W każdej konfiguracja miała co najmniej
sekundę rozgrzewki i 20 prób. Wyniki przedstawiono jako średnie
z 60 prób; nie odrzucano wolniejszych pomiarów.

Czas mierzono zegarem monotonicznym biblioteki standardowej C++.
Dla CPU równoległego uwzględniono tworzenie i zakończenie wątków.
Dla GPU pomiar obejmował rezerwację pamięci, transfery, obliczenia,
synchronizację i zwolnienie pamięci. Generowanie danych i kontrola
wyników odbywały się poza pomiarem. Wszystkie 1440 prób głównej
serii dało poprawny wynik.

Przyspieszenie względem wersji sekwencyjnej obliczono jako
\begin{equation}
S=\frac{\overline{t}_{\mathrm{seq}}}{\overline{t}_{\mathrm{wariant}}}.
\end{equation}
Na wykresach punkty oznaczają średnie, a przerywane linie
dopasowany trend potęgowy. Słupki na wykresach czasu pokazują
jedno odchylenie standardowe. Punkty nie są połączone.
Wyniki zaokrąglono do dwóch miejsc po przecinku po wykonaniu obliczeń.

GPU profilowano osobno narzędziem NVIDIA Nsight Systems dla $N=512$.
Analiza dotyczyła 20 prób po rozgrzewce. Czasów z tego uruchomienia
nie używano do obliczania przyspieszeń ani głównych wykresów.

\section{Wyniki i wnioski}

\subsection{Czas wykonania i przyspieszenie}

Tabela~\ref{tab:czasy} i rysunek~\ref{fig:czasy} porównują trzy
wersje algorytmu. W tym zestawieniu CPU równoległy używa 8 wątków.

\begin{table}[H]
\caption{Średni czas mnożenia [ms]}
\label{tab:czasy}
\centering
\begin{tabular}{@{}rccc@{}}
\toprule
$N$ & CPU sekw. & CPU, 8 wątków & GPU\\
\midrule
64 & 0,15 & 0,21 & 0,14\\
128 & 2,33 & 0,67 & 0,16\\
256 & 18,41 & 3,23 & 0,29\\
512 & 385,62 & 61,18 & 0,90\\
\bottomrule
\end{tabular}
\end{table}

\wykres{../results/2026-11-01/time.pdf}
{Czas wykonania w zależności od rozmiaru macierzy.}
{fig:czasy}

Wraz ze wzrostem macierzy zwiększał się czas wszystkich wersji.
Dla $N=512$ GPU potrzebowało 0,90\,ms, CPU z 8 wątkami
61,18\,ms, a CPU sekwencyjny 385,62\,ms.
Odpowiada to przyspieszeniu 428,32 razy dla GPU i 6,30 razy dla CPU
równoległego, co przedstawia rysunek~\ref{fig:speedup}.
Wyniki odnoszą się do prostej implementacji i podanych warunków badań.

\wykres{../results/2026-11-01/speedup.pdf}
{Przyspieszenie względem CPU sekwencyjnego.}
{fig:speedup}

Dla $N=64$ CPU z 8 wątkami było wolniejsze od wersji sekwencyjnej,
a przewaga GPU była niewielka. Pokazuje to, że równoległość daje
większe korzyści, gdy obliczeń jest wystarczająco dużo.

\subsection{Skalowalność CPU i zmienność pomiarów}

Rysunek~\ref{fig:watki} pokazuje wpływ liczby wątków.
Dla $N=512$ zwiększenie ich liczby z 1 do 8 skróciło czas
z 386,06 do 61,18\,ms. Dla $N=64$ najlepsze były 4 wątki
z czasem 0,12\,ms, krótszym także od GPU. Przy małych zadaniach
koszt obsługi kolejnych wątków ogranicza skalowanie.

\wykres{../results/2026-11-01/cpu_scaling.pdf}
{Czas wykonania CPU równoległego w zależności od liczby wątków.}
{fig:watki}

Dodatkowe próby dla $N=511$, 512 i 513 potwierdziły spowolnienie
CPU dla $N=512$. Może ono wynikać z kolumnowego odczytu macierzy $B$
i działania pamięci podręcznej; nie potwierdzono jednak tej przyczyny
pomiarem chybień. Złożoność $O(N^3)$ opisuje liczbę operacji,
a nie stały czas każdej z nich.

Ustawienia zegarów nie usunęły całej zmienności. Dla GPU przy
$N=512$ odchylenie standardowe wyniosło 0,13\,ms, czyli 14,40\%
średniej. Pełny czas obejmuje też sterownik, transfery i obsługę
pamięci. Nie ustalono osobnego wpływu tych czynności na rozrzut.

\clearpage
\subsection{Profilowanie GPU}

Tabela~\ref{tab:profil} zawiera wybrane wyniki Nsight Systems.
Czasy opisują pracę urządzenia, a nie całe wywołanie programu.

\begin{table}[H]
\caption{Profil GPU dla $N=512$: średni czas z 20 prób [ms]}
\label{tab:profil}
\centering
\begin{tabular}{@{}lr@{}}
\toprule
Czynność & Czas\\
\midrule
Obliczenia na GPU & 0,45\\
Przesłanie obu macierzy do GPU & 0,09\\
Odebranie wyniku z GPU & 0,05\\
\bottomrule
\end{tabular}
\end{table}

Jądro zajmowało 76,33\% łącznego czasu obliczeń i transferów
na urządzeniu. Profil wykazał też koszt powtarzanej rezerwacji
i zwalniania pamięci. Ponowne użycie buforów mogłoby go ograniczyć,
ale takiej optymalizacji nie wprowadzono.

\section{Podsumowanie}

Przygotowano i sprawdzono trzy wersje mnożenia macierzy.
Dla większych badanych rozmiarów GPU było najszybsze,
a wielowątkowość skróciła czas CPU. Dla małych macierzy
koszt uruchamiania pracy równoległej ograniczał korzyści.
Pomiary, analiza skalowalności i profilowanie pozwoliły
osiągnąć cel projektu.

\begin{thebibliography}{2}
\raggedright
\bibitem{nvsmi} NVIDIA, \emph{Dokumentacja nvidia-smi},
\url{https://docs.nvidia.com/deploy/nvidia-smi/index.html},
dostęp: 01.11.2026.
\bibitem{cpufreq} The Linux Kernel, \emph{CPU Performance Scaling},
\url{https://docs.kernel.org/admin-guide/pm/cpufreq.html},
dostęp: 01.11.2026.
\end{thebibliography}
\end{document}
